\section{AegisKV Design}
\label{sec:design}

AegisKV manages KV across three protection tiers: selective Peer HBM,
shared DRAM, and shared storage. Peer HBM holds extra hot-block replicas
on a preassigned takeover DP; each admitted DRAM checkpoint is complete.
Under pressure, a complete checkpoint can move from DRAM to shared
storage. Ordinary execution HBM and transient recovery workspace are
accounted separately from the Peer budget. The protected domain is a
fail-stop compute worker: the directory, token log, takeover target and
advertised backing must survive that failure. Same-host memory or storage
does not protect against loss of that host. This section specifies the
proposed design; constructed figures are not controller measurements.

\subsection{Design Insights}
\label{sec:design-insights}

\paragraph{Prioritize workflow exposure, while preserving normal reuse.}
A faster restore advances a workflow only when its delay exceeds the
available dependency slack (Figure~\ref{fig:motivation-recovery-exposure}).
However, a large-slack agent can reuse its history after a tool call.
Assigning that agent permanently to recomputation would waste serving
capacity. AegisKV therefore separates \emph{preserving a reusable copy}
from \emph{retaining it in the faster tier}. Shared immutable objects
serve both purposes without duplicating their physical byte charge.

\paragraph{Sparse access identifies candidates, not a complete backup.}
Native sparse attention can use only a subset of historical KV for a
query, but future selections remain uncertain. Figure~\ref{fig:design-insights}
shows that causal access history predicts future block demand unevenly
across models. At a 40\% block budget, past-frequency coverage is 77.8\%
for NOSA, 53.6\% for MiniCPM4, and 43.0\% for the experimental Seer path.
The corresponding same-budget oracle means are 93.7\%, 59.0\%, and
43.3\%. Seer's low absolute coverage is therefore largely a budget limit
for this future-query union, rather than a large gap to its oracle.
These observations motivate selective fast-tier retention and a mandatory
miss path, rather than discarding the remaining checkpoint. The three
documents per model support this observed contrast, not a population-wide
prediction guarantee or a measured recovery speedup.

\begin{figure}[t]
 \centering
 \includegraphics[width=\columnwidth]{figures/retry-kv-prediction.pdf}
 \caption{Historical access predicts distinct future KV unevenly.
 Three native PG19 captures/model use 32K history: queries
 32,257--32,768 predict blocks needed by 32,769--32,832.
 Objects are layer/64-token blocks containing all KV heads;
 mandatory blocks count against the 40\% budget. Seer uses an
 experimental Qwen3-8B Prefill gate~\cite{song2026seerqwen3}, threshold
 $10^{-3}$, with final-block dense override disabled.
 Bars/whiskers are document means/ranges, not recovery runs. Diamonds
 show offline oracle mean coverage with the same mandatory blocks and
 capacity; future demand is unavailable to the online predictor.}
 \label{fig:design-insights}
 \Description{Held-out native block coverage varies across three models;
 weak effects remain visible.}
\end{figure}

\paragraph{Overlap determines the marginal value of Peer memory.}
A complete DRAM checkpoint avoids historical reconstruction but still
requires H2D transfer. Peer-resident blocks on the takeover DP can remove
that transfer. Early-layer demand may delay replay startup; later demand
can overlap with computation that still has to execute. The useful
resident depth therefore depends on checkpoint length, uncommitted tail,
demand reveal time and resource contention. The Master chooses per-request
layer depth and hot sets; a separate Adapter adjusts only the total Peer
capacity available to those choices.

\subsection{System Overview}
\label{sec:system-overview}

Figure~\ref{fig:system-overview} retains three ownership layers.
The \emph{Application} exposes agent requests, dependencies and progress.
The \emph{Controller} contains the Workflow Tracker, Recovery-aware
Router, KV Master and surviving directory. The Tracker supplies duration
and slack estimates. At arrival, the Router selects an execution DP and
binds a failure-separated takeover DP, which is also the Peer replica
destination. The Master jointly values whole-checkpoint DRAM/storage
placement and optional Peer hot replicas. Its current-value comparisons
include victims, workflow dependencies and transition costs.

The \emph{Engine} contains symmetric serving DPs, each with ordinary KV
and a bounded Peer region in HBM. Each DP's Adapter supplies a window's
Peer capacity budget; its Connector executes versioned transfers and
layered recovery. Shared DRAM and storage lie outside the protected DP
boundaries. The directory records manifests, copy generations, leases
and target bindings; its metadata and ordered token log must survive.

\begin{figure}[t]
 \centering
 \includegraphics[width=\columnwidth]{figures/system-overview.pdf}
 \caption{Three-layer architecture and three KV protection tiers.
 The Router binds the takeover/Peer DP before placement. KV Master
 chooses per-request L/K and complete DRAM/storage (D/S) placement;
 each Adapter budgets Peer capacity, and Connectors transfer KV.
 Blue links carry KV, dark dashed links carry control or metadata,
 and red dashes show retry. Two symmetric DP boundaries contain
 Local KV and Peer KV; complete shared backing lies outside both.
 GPU/HBM denotes a DP's compatible GPU/TP group.}
 \label{fig:system-overview}
 \Description{Application, Controller and Engine bands contain a workflow,
 Tracker, Router, Master, directory, symmetric DP adapters/connectors,
 partial Peer replicas and full shared checkpoints.}
\end{figure}

A physical object contains one layer's KV shard for a sealed token block.
Identity includes model/weight revision, prefix and position semantics,
layer, TP shard, format, and token interval. A legal placement group
includes the objects and selector state required by the native engine;
logical sparse selections, physical storage blocks, and kernel tiles need
not coincide. The engine chooses the physical block size; the 64-token
unit in Figure~\ref{fig:design-insights} is an experimental convention.

Prefill offers newly sealed groups. During decode, a configured,
block-aligned interval $Z$ triggers backup review. Arrivals, tool handoffs,
material demand/slack changes, and capacity pressure also trigger review.
A decision does not imply immediate commitment: if $M$ is the last
complete checkpoint frontier, the replay tail $N$ may exceed $Z$ when
copies lag. A complete checkpoint includes all historical dependencies needed to
resume native attention: every layer and TP shard, selector state and the
matching log frontier. A manifest advertised as DRAM-resident resolves all
of these dependencies in DRAM; it is not a hot subset assembled across
tiers. Storage holds complete lower-tier checkpoints. Shared immutable
objects and incremental writes can avoid duplicating unchanged bytes,
but a new manifest becomes complete only when all its dependencies exist.
During whole-checkpoint migration, the old complete location remains
authoritative until the new one commits.

\subsection{Workflow Tracker: Remaining Time and Slack}
\label{sec:workflow-protection}

The Tracker builds a residual DAG of unfinished execution segments.
Fork/join boundaries split parent work so child execution is not counted
again as parent waiting. A gradient-boosted decision-tree model (GBDT)
predicts remaining duration. Its inputs group three types of causal
information: role/model and current input length; elapsed execution and
completed-call count; and completed calls' output lengths, API durations,
context growth, and inter-call gaps. Tool type and task descriptors can
be added when the orchestrator exposes them. Current-call output length,
future tool results, and future call count are unavailable inputs.
Entry predictions mask missing history; elapsed time is not subtracted
again from an already remaining-time prediction.

Figure~\ref{fig:tracker-correlations} examines six observable numerical
inputs in the CC/Weka trace~\cite{semianalysis2026cctraces}. Each of 1,443
eligible agents contributes one snapshot at its second request, after
its first call completes. The target is the remaining duration of the
whole agent, including later calls and gaps. Input-length features are
64-token hash-prefix proxies; explicit task semantics and tool outcomes
are not recorded. All six marginal correlations are weak, and their
workflow-bootstrap intervals cross zero. Thus, a single-feature rule
is not established by this cohort. These associations neither measure
GBDT feature importance nor validate its nonlinear predictions or the
residual-segment target required by the online Tracker.

\begin{figure}[t]
 \centering
 \includegraphics[width=\columnwidth]{figures/tracker-correlations.pdf}
 \caption{Tracker evidence at request 2: 1,443 agents in 170 workflows,
 one snapshot per agent. (a) Spearman associations with whole-agent
 remaining time; bars are percentile 95\% intervals from 2,000
 whole-workflow bootstrap samples. Asterisks mark prefix-block length
 proxies. (b) Empirical absolute-error CDFs from existing five-fold,
 workflow-disjoint predictions: training-fold median, basic GBDT, and
 GBDT with call history. All tails are retained; the horizontal scale is
 logarithmic above 1\,s. Neither panel validates residual-segment Slack.}
 \label{fig:tracker-correlations}
 \Description{Six weak feature associations and whole-agent held-out
 error distributions. GBDT errors remain substantial; history does not
 uniformly improve the basic model.}
\end{figure}

Figure~\ref{fig:tracker-correlations}(b) makes prediction error explicit.
The existing workflow-separated five-fold pilot gives mean absolute
errors of 438.6, 406.7, and 407.1\,s for median, basic GBDT, and
history GBDT. Their median absolute percentage errors are 76.2\%,
67.1\%, and 69.3\%. The basic-minus-median MAE difference has a paired
workflow-bootstrap 95\% interval of $[-80.9,4.9]$\,s; it does not
establish a reliable mean-error reduction. Fitted predictions are fixed
in this bootstrap. Only 36.5\% of basic-GBDT predictions are within 50\%
of actual remaining time. The cohort comprises completed subagents and
includes future inter-call gaps, so it cannot validate runtime segment
predictions. Deployment validation must separate whole workflows and
time periods and report false-positive slack, not just duration error.
Published trajectory
prediction~\cite{zhang2026heddle} motivates feature construction but does
not establish wall-clock accuracy for AegisKV.

Critical-path analysis~\cite{kelley1959cpm} computes earliest and latest
finishes; their nonnegative difference is slack $\sigma_r$ inherited by a
request's segment. Requests on the same segment do not each receive an
independent copy of that margin. A refresh costs $O(|V|+|E|)$ after
prediction. Missing dependencies or invalid predictions fall back to
request recovery cost and normal reuse with zero slack. A new DAG version
invalidates outstanding valuation snapshots.
Without validated segment forecasts, the policy uses the zero-slack
fallback rather than treating the pilot's point estimates as safe
demotion margins. The current Tracker evidence motivates this guard;
it does not establish predictive accuracy as a system contribution.



\subsection{Recovery-aware Router}
\label{sec:recovery-router}

Normal dispatch balances load and compatible prefix locality, retaining
Agent affinity when useful. The CC/Weka trace~\cite{semianalysis2026cctraces}
shows 97.90\% potential same-agent token coverage, with 0.37\% additional
coverage from other agents; these are potential prefixes, not cache hits. A tool handoff drops
an execution reference and the request's logical Peer pin. The complete
cached history remains eligible for reuse; tool waiting is not deletion.

\begin{figure}[t]
 \centering
 \includegraphics[width=\columnwidth]{figures/agent-prefix-reuse.pdf}
 \caption{Agent-local prefix reuse in the CC/Weka trace:
 98,827 requests, 2,090 agents, and 393 workflows. Coverage uses earlier
 completed requests with the same recorded model and workflow, measured
 by prefix-block hashes. Same-agent and other-agent coverage overlap;
 only the additional cross-agent portion extends same-agent coverage.
 These are potential token-prefix reuse rates, not runtime cache hits.}
 \label{fig:agent-prefix-reuse}
 \Description{Same-agent history covers most observed input tokens;
 other-agent history adds little overall but helps first subagent calls.}
\end{figure}

Figure~\ref{fig:agent-prefix-reuse} separates overlapping coverage from
additional coverage. Cross-agent reuse can still matter on a subagent's
first call; affinity is therefore a preference, not a prohibition on
sharing or load balancing.

At arrival, the Router records a versioned takeover binding in a separate
failure domain. The Master places Peer objects on that target's matching
TP shards. Source overload or target pressure can cause an explicit
rebind: the old target remains authoritative until the new binding and
its advertised copies are valid. A rebind need not wait for a Peer copy;
it can publish the complete DRAM/storage path with Peer coverage zero.

On a confirmed failure, the Router fences the old source epoch, resolves
the last complete manifest and activates the bound healthy target with
feasible workspace. If that target is unavailable, it chooses another
compatible survivor using the complete backing path; copies on an
unselected Peer are not treated as local HBM hits. Queuing and readmission remain part of
recovery cost. Before acknowledging output, a surviving ordered log
records request ID, token position/value, and source epoch; prompt and
accepted tool inputs are also recoverable. Replay suppresses duplicate
delivery and rejects stale callbacks. Exactly-once external tool side
effects remain the orchestrator's responsibility.

\subsection{KV Master: Three-Tier Placement and Direct Revaluation}
\label{sec:kv-master}

\subsubsection{Full-checkpoint placement and Peer admission}
\label{sec:placement-value}

The Master compares complete, achievable placement plans. For an isolated
interrupted segment, recovery time $R_r$ exposes delay
$\max(0,R_r-\sigma_r)$. A common worker failure can interrupt several
branches: if two required branches each finish 10\,s late, accelerating
only one to 2\,s leaves their join 10\,s late. Thus, isolated block scores
cannot be added to obtain workflow benefit.

For current plan $P$ and candidate $P'$, we evaluate
\begin{equation}
 \Delta J=A(P',P)+\sum_e p_e\sum_w
       \bigl[T_w(P',e)-T_w(P,e)\bigr].
 \label{eq:value}
\end{equation}
Episodes $e$ cover normal next use and named compute-domain failures over
a fixed horizon. The weights $p_e$ are declared probabilities or expected
episode counts, not raw failure rates. A failure includes every request
interrupted in that domain. Profiles estimate additional segment delays
from surviving coverage, tail replay, I/O and target queues; a max-plus
DAG pass yields $T_w$. The transition term $A$ accounts for foreground
delay from copying and temporary occupancy, excluding costs already in
the episode terms. All terms have units of time.

Normal reuse remains in this objective even when failure probability is
zero or slack is large. The Master distinguishes three decisions.
Complete checkpoints compete for DRAM against a complete storage-backed
continuation. Peer hot replicas compete against the \emph{current full
backing path}, including the H2D wait still exposed after overlap.
Last-copy deletion is valued against a remaining checkpoint or log-based
reconstruction. Peer is optional acceleration, not the sole checkpoint:
a committed Peer replica always has a committed complete backing manifest
in DRAM or storage. A checkpoint may remain complete in storage while
selected blocks stay on the Peer. Its slower miss path is included in
revaluation; demotion does not silently preserve the old DRAM cost. Every shared owner is considered,
while each physical replica consumes capacity once. A plan may forecast future commitment for cost estimation, but pending
copies never count as current recoverable coverage.
No future failure time, prompt, or query is available to the policy.

For newly sealed state, the Master proposes a complete checkpoint extension
in DRAM; if it does not fit, it considers direct commitment to storage.
Only new objects need transfer when immutable dependencies already exist.
If neither tier can accept a complete extension, the old frontier stays
valid and the replay tail grows. A large slack does not disable this
normal-cache policy, and finite storage can still force an explicit
reconstruction fallback. After backing placement, the Master considers
optional Peer admission and changes to existing hot replicas. This is
not a permanent request-level choice among recompute, DRAM and Peer.

\paragraph{Why a priority label is insufficient.}
Slack measures tolerated delay, not how much delay a placement can remove.
Figure~\ref{fig:retention-value} applies the isolated-fault exposure rule
to two declared recovery profiles. At zero slack, retaining A saves
2\,s using two capacity units; at 6\,s slack, retaining B saves 10\,s
using one. Sorting only by slack reverses their value-per-byte order.
Recency can likewise miss a valuable checkpoint with no recent normal
access, while raw saved time ignores how many bytes must remain resident.
These observations require cost- and size-aware comparison. They do not
by themselves select a particular replacement algorithm.

\begin{figure}[t]
 \centering
 \includegraphics[width=\columnwidth]{figures/retention-value.pdf}
 \caption{An analytical counterexample to slack-only priority.
 Per-unit saving is $([R_{\mathrm{slow}}-\sigma]^+
 -[R_{\mathrm{fast}}-\sigma]^+)/b$, conditional on an isolated failure.
 A occupies two normalized capacity units with restore times 4/2\,s;
 B occupies one with times 18/8\,s. Marked points compare A at slack 0
 and B at slack 6. Normal-use value is set to zero to isolate the effect.
 Inputs are constructed, not hardware measurements or a guarantee of
 optimal placement.}
 \label{fig:retention-value}
 \Description{A zero-slack group can have much lower saved workflow
 time per unit of capacity than a group with positive slack.}
\end{figure}

\subsubsection{Choosing per-request layers and hot blocks}
\label{sec:sparse-placement}

For request $r$, $L_r$ is the number of early layers eligible for Peer
residency, and $K_{r,l}$ is the number of engine token blocks replicated
in layer $l\leq L_r$. Each block includes all required TP-shard objects. Each such layer has
a selected hot set; the complete
checkpoint remains in DRAM or storage. The engine's native selections
are unchanged. A bounded causal access summary ranks blocks and includes
mandatory sink/window dependencies and the selector state needed by the
candidate. Logical selection units need not equal physical blocks.

The Master enumerates a configured, bounded set of prefix depths and
per-layer hot-set budgets, including $L_r=0$ and the current placement.
For each candidate it schedules predicted demand, misses, H2D copies and
remaining layer computation with finite workspace. It estimates the
change in Equation~\ref{eq:value}, charges newly allocated physical bytes
and movement, and compares the candidate together with any victims.
Small bundles across jointly delayed branches are also candidates within
the same work budget. The resulting $L_r$ and $K_{r,l}$ are per-request
decisions, not an Adapter-wide quota or independent block-score sums.
Bounded enumeration can miss a better plan; no optimality or
submodularity claim follows.

\begin{figure}[t]
 \centering
 \includegraphics[width=\columnwidth]{figures/sparse-recovery.pdf}
 \caption{Complete backing and a partial Peer replica. All historical
 B0--B3 blocks in three layers reside in DRAM; the dashed tail is
 uncommitted. Peer holds B0/B3 in the first two layers ($L=2$, $K=2$).
 Layer 1's native demand for B0/B2 triggers a DRAM fetch for B2;
 B3 was an unused prediction. Whole-checkpoint demotion can replace
 the full DRAM backing with complete shared storage. The example
 describes state, not measured prediction or recovery accuracy.}
 \label{fig:sparse-recovery}
 \Description{Full DRAM history, an uncommitted tail and selected Peer
 replicas; a prediction miss fetches a block from the complete backing.}
\end{figure}

Context growth does not force either parameter to grow at every token.
An engine with a fixed sparse budget can retain the same $K_{r,l}$ while
block identities change; a proportional budget can increase it. Even in
the latter case, Peer admission can cap $K$ under finite capacity.
Reconsideration occurs on the configured $Z$ review and material events,
not on every decode token. A proposed replacement set is distinguished
from the committed set until transfers finish. The profile uses current
checkpoint length $M$ and a declared distribution or bound for unknown
tail $N$, with support for copy lag beyond $Z$. It cannot use the future
failure phase or infer a universal monotone relation between $Z$ and $L$.

Figures~\ref{fig:design-pipeline} and~\ref{fig:overlap-frontier} isolate
Peer's incremental effect over the same complete DRAM checkpoint.
Only computation still required after takeover can hide transfers;
avoided historical prefill is not overlap time. Early residency can
remove startup waits, while later copies may already be hidden. If the
profile predicts no useful exposed wait reduction, $L_r=0$ is a valid
choice even for a zero-slack request.

\begin{figure}[t]
 \centering
 \includegraphics[width=\columnwidth]{figures/recovery-pipeline.pdf}
 \caption{Matched analytical Peer comparison: four equal-sized demand
 groups, H2D 4\,ms/group and replay 3\,ms/layer. The same complete
 checkpoint is in DRAM; all demand is correctly revealed at time zero.
 Separate transient workspace holds four groups. DRAM-only recovery
 takes 19\,ms; equal one-group Peer budgets at layer 4 and layer 1
 yield 18 and 15\,ms. Bar labels are layer numbers. Late-layer residency
 is a counterfactual for position, not the prefix policy. Rates are
 constructed; transfer, compute and admission queues are not measured.}
 \label{fig:design-pipeline}
 \Description{Equal-byte early and late Peer copies remove different
 portions of the same DRAM-to-HBM recovery wait.}
\end{figure}

\begin{figure}[t]
 \centering
 \includegraphics[width=\columnwidth]{figures/overlap-frontier.pdf}
 \caption{Per-request Peer depth under three fixed replay profiles.
 Each layer contributes one equal-sized required group ($K=1$), with
 the same full DRAM checkpoint and H2D cost as
 Figure~\ref{fig:design-pipeline}. Each $L=0$ is its own matched
 DRAM baseline. Flat regions show transfers already hidden by replay;
 larger compute cost can reduce useful depth while increasing latency.
 These fixed-tail analytical cases are not a distribution of failures.}
 \label{fig:overlap-frontier}
 \Description{Recovery versus Peer prefix depth saturates at different
 depths for different remaining computation.}
\end{figure}

\subsubsection{Online admission, direct revaluation, and eviction}
\label{sec:placement-plan}

\paragraph{Compare current plans, not historical credits.}
The Master re-evaluates placement at the events described in
Section~\ref{sec:system-overview}. Each decision uses one versioned snapshot
of the DAG, demand summaries, committed copies, queues, leases, and byte
ledgers. A cached estimate is reusable only while these dependencies are
unchanged. There is no rent level, accumulated credit, or renewal rule.
A running request logically pins its committed Peer hot replicas against
replacement by another request. Request completion or tool handoff releases
that pin and makes those replicas eligible for eviction; the complete
backing history remains reusable. Normal reuse and failure exposure
both remain in Equation~\ref{eq:value}.

Direct revaluation is a design choice, not a claim of optimal search.
The isolated comparison in Section~\ref{sec:eval-rental} retains the
rejected rental alternative: historical credit screening increased
modeled miss-plus-movement cost without reducing replacement-value checks.
That additive example does not validate the complete DAG controller or
measure its CPU cost. The runtime therefore bounds candidate and victim
evaluations and reports both decision work and resulting placement regret.

\paragraph{Construct a feasible replacement.}
Candidates have tier-specific granularity: a complete checkpoint for
DRAM or storage placement and a legal per-layer block bundle for Peer.
Shared physical dependencies count once; evicting one logical checkpoint
releases only objects that no other retained manifest needs. The safe
victim set $Q$ excludes physical read leases and pending copy dependencies.
Ordinary Peer competition also excludes other running requests' logical
pins. Self-replacement of a request's predicted hot set is allowed, and
Adapter pressure may override a logical pin after safe redirection. Demotion needs a committed lower copy or a
feasible copy-and-commit schedule; last-copy deletion needs an explicit
valid fallback. Logical Peer pins never authorize freeing memory under a physical read lease. Shared bytes are charged once, while
the plan's effects on all owners are evaluated together. A victim that
releases zero physical bytes cannot provide capacity relief on its own;
it is skipped unless a legal owner bundle releases shared bytes.

For each candidate, the Master constructs a private victim plan. It scans
a bounded, rotating set of eligible checkpoint manifests or Peer bundles in the affected tier, ordered
by stable identity from a persistent cursor. Starting from the current
placement, it computes each feasible removal's additional cost through
Equation~\ref{eq:value} per byte released in that tier. It appends the
least-loss removal, breaking ties by fewer moved bytes and then identity,
and re-evaluates subsequent removals against the changed tentative plan.
This continues until sufficient room exists for the incoming group and
its transfer workspace, or the configured work budget is exhausted.
The cursor resumes at the next eligible ID on later decisions; a bounded
scan can miss better victims and supplies no global optimality guarantee.

A whole-checkpoint DRAM demotion that lacks storage space first constructs a lower-tier
relief plan by the same rule, with storage deletion as the final tier.
It includes those victims, fallback changes, and copy dependencies in
the same plan. It never assumes an infinite lower tier. A backing demotion also revalues every dependent Peer plan under its new
miss path. A tentative
removal is eligible only if its transition can be scheduled within the
physical and reserved byte limits; incoming copies wait for real release.
If no such plan fits the work budget, this candidate is skipped. The
ratio only constructs the victim set: the final score is the complete
candidate-plus-victim cost, not the sum of independent block scores.

\begin{figure}[t]
 \centering
 \includegraphics[width=\columnwidth]{figures/placement-competition.pdf}
 \caption{Direct replacement comparison in a constructed additive
 example with equal-size Peer bundles, valid full backing and cost units
 of seconds. Candidates y/z have no running logical pin. Promoting x
 saves 6; releasing y or z costs 2 or 5; either transition costs 1.
 Thus the complete changes are $1+2-6=-3$ and $1+5-6=0$.
 A leased replica is excluded even if its removal cost is lower.
 With a margin below 3, the first plan is eligible for reservation;
 physical release still waits for valid backing and reader drainage. General
 common-failure plans use the joint DAG, not this additive shortcut.}
 \label{fig:placement-competition}
 \Description{Current values compare two complete candidate and victim
 plans; a leased object is excluded and only a profitable safe plan proceeds.}
\end{figure}

\paragraph{Choose and apply one plan.}
Among the fully evaluated plans, the Master selects the smallest
$\Delta J$, breaking ties by fewer moved bytes and stable object IDs.
An optional plan must improve predicted cost by more than the configured
uncertainty allowance plus migration hysteresis. These allowances are
heuristic guards unless separately calibrated; they are not statistical
confidence bounds. Rejected plans change neither placement nor capacity.
Before acceptance, the Master rechecks snapshot versions and reserves
the immediately executable destinations and transfer workspace. It queues
dependent actions in copy--commit--redirect--drain--release order and
starts incoming copies only when their space is physically available.
Later completions trigger a new snapshot; tentative plans from an older
placement are not reused as if nothing had changed.

Required capacity relief has no incoming bidder. The same current-loss
rule constructs a safe relief plan until the target is met, re-evaluating
after each tentative removal. It may accept a positive cost to satisfy
the capacity constraint. Storage deletion considers unowned or superseded
state first. Further relief may withdraw an active checkpoint claim,
but must first publish a complete older manifest or a prompt-and-token-log
fallback that does not depend on the deleted object. All owners' costs
enter the plan, and old-generation readers must drain before deletion.
If no executable plan exists within the budget, new admission pauses;
the old frontier remains valid and a later event resumes the search.

\par\medskip\noindent\begin{minipage}{\columnwidth}
\raggedright
 \refstepcounter{designalgorithm}\label{alg:placement}
 \textbf{Algorithm \thedesignalgorithm: Direct placement revaluation}
 \par\smallskip
 \begin{tabularx}{\columnwidth}{@{}r@{\quad}X@{}}
 \toprule
 1 & On a review event, snapshot versions, copies, queues and byte ledgers \\
 2 & Refresh M/N, causal demand and affected values; reuse existing objects \\
 3 & Read Adapter budgets; on shrink, construct safe current-loss relief \\
 4 & Otherwise, form bounded complete-checkpoint and Peer L/K candidates; best $\leftarrow$ none \\
 5 & For each candidate within the work budget: \\
 6 & \quad Construct safe victims by re-evaluated loss per released byte \\
 7 & \quad Include lower-tier victims and fallback changes; skip if infeasible \\
 8 & \quad Evaluate the complete plan with Equation~\ref{eq:value} \\
 9 & \quad Keep the lowest-cost plan that clears uncertainty and hysteresis \\
 10 & If none qualifies, retain backing; on pressure, pause new admission \\
 11 & Recheck versions; on mismatch, reschedule from a fresh snapshot \\
 12 & Reserve executable space and enqueue dependent transitions \\
 13 & On completion/failure, reconcile actual state and trigger a fresh review \\
 \bottomrule
 \end{tabularx}
 \par\smallskip\noindent\textit{Predicted release is not free memory; a plan is not a committed copy.}
 \end{minipage}\par\medskip

Profiles carry model/layout, size, tail, topology, queue load, support
range, uncertainty, and expiry. Missing or stale profiles freeze optional
promotions and use bounded LRU for required relief under the same safety
protocol. Candidate construction, each victim comparison, transfer
profiling, and complete DAG passes count against the work budget; report
their counts, decision latency, and the loss from bounded search. The
isolated model's replacement checks are not measurements of these costs.

\subsection{Adapter: Windowed Peer Capacity}
\label{sec:adapter}

Each DP's Adapter measures ordinary execution demand, reserved workspace
and recent pressure over a time window. It sets the next window's Peer
budget to the smaller of a configured cap and the remaining managed HBM
pool after the serving allocation and a headroom reserve, clipped at zero.
The serving allocation includes live execution state and near-term demand;
it is not the Peer allocation subtracted a second time. The cap can be a
configured fraction of the managed pool. No fixed window duration or
failure-rate-controlled checkpoint interval is required: $Z$ remains a
system setting. The Adapter does not choose $L_r$, $K_{r,l}$ or hot blocks.

A lower target becomes a relief request to the Master. New Peer admission
pauses while committed plus reserved occupancy exceeds the target. The
Master can retire running requests' logically pinned Peer replicas in
order of current whole-plan loss; their complete backing remains valid.
Physical readers and transfers must drain before the bytes become free.
A valid immutable lower copy means no redundant D2H write is needed.
If backing cannot be validated, release must wait for repair or an explicit
valid fallback; lowering a budget never manufactures free memory.
Emergency foreground pressure may request relief within a window under
the same protocol. If leases prevent immediate relief, the engine applies
backpressure rather than reusing occupied addresses.

\subsection{Connector: Committed Copies and Layered Recovery}
\label{sec:connector}

Connectors implement placement plans under byte reservations and bounded
queues for foreground reuse, demanded recovery and checkpoint movement.
Configured service shares and bounded transfer quanta limit interference;
demanded misses precede speculative reads within the recovery share.
A request's newly sealed objects extend complete backing first. Only
objects covered by a committed complete manifest are advertised as Peer
backup; an in-flight accelerator copy cannot advance that manifest.
This design trades possible commit lag for an unambiguous safe lower path.
Copying and Peer occupancy can still degrade healthy serving.

\begin{figure}[t]
 \centering
 \includegraphics[width=\columnwidth]{figures/checkpoint-lifecycle.pdf}
 \caption{Three-tier commitment and release. (a) Pending tokens
 129--192 leave $M=128$ unchanged until all dependencies commit.
 (b) A complete DRAM checkpoint is freed only after complete storage
 commitment, redirection and reader drainage. (c) A Peer replica has
 complete backing; releasing it requires no new D2H copy when that
 backing is valid. Numbers illustrate ordering, not fault-injection results.}
 \label{fig:checkpoint-lifecycle}
 \Description{A complete manifest frontier differs from an in-flight
 copy; storage commitment and reader drainage precede DRAM release; Peer release preserves a complete lower path.}
\end{figure}

Objects move through pending, committed, and retiring generations
(Figure~\ref{fig:checkpoint-lifecycle}). Advancing a frontier requires every
layer, TP shard, selector state, token interval, and matching surviving
log prefix. It is not the largest individually copied token position.
DRAM demotion follows copy--verify--commit--redirect--drain--free. The
backend must meet the claimed fault-domain persistence contract; buffered
write completion alone does not imply durable storage. Copy failure keeps
the previous authoritative location and releases only failed reservations.
Duplicate or stale-epoch callbacks cannot advance coverage.

Storage last-copy deletion separately publishes a complete older
checkpoint or prompt-and-token-log fallback before draining readers.
Any Peer replica no longer covered by a complete backing manifest is
withdrawn and retired as part of that same transition. A remaining hot
subset is never advertised as a complete recovery checkpoint.
Whole-checkpoint migration may temporarily duplicate objects, and these
bytes count against actual destination capacity; partial transfer never
makes the source checkpoint reclaimable.

Recovery acquires a manifest, target binding and location leases.
Compatible hot replicas already on the takeover DP bypass H2D. Objects
adopted into execution are reclassified from Peer to execution accounting
without charging the physical bytes twice or releasing their addresses.
Other
blocks come from the complete DRAM checkpoint, or are read from shared
storage before H2D. The transfer path prefetches predicted native groups in layer order into bounded transient
workspace, while the engine replays the acknowledged tail. On-demand
selection resolves misses before executing the affected operation. Dense
layers use full demand; unavailable sparse metadata is restored or
reconstructed before selection. TP ranks obey matching collectives, and
batched replay is used only when it preserves the engine's native
semantics. A failure during prefill also starts from the last complete frontier,
possibly zero. Speculation may improve timing but never defines correctness.
